\section{Построение решения и исследование задачи}
\subsection{Выбор подходящих инструментов}

Из-за огромных и постоянно растущих объёмов трафика разрабатываемая система должна обладать высокой пропускной способностью, низкой и предсказуемой задержкой, эффективным использованием многопроцессорных архитектур и минимальными накладными расходами на управление памятью.
Языки программирования с динамической типизацией и автоматическим контролем памяти (сборкой мусора) не подходят для данных требований, поскольку недетерминированные паузы GC, overhead рантайма и отсутствие контроля над расположением данных в памяти приводят к высоким накладным расходам и деградации производительности при пиковых нагрузках \cite{realtime2021gc}.

% Для формализации требований к системе были определены целевые метрики производительности, представленные в таблице~\ref{tab:performance_targets}.
% \begin{table}[h]
% \centering
% \caption{Целевые метрики производительности системы событий}
% \label{tab:performance_targets}
% \begin{tabular}{@{}lll@{}}
% \toprule
% \textbf{Метрика} & \textbf{Целевое значение} & \textbf{Обоснование} \
% \midrule
% Пропускная способность &
% ≥ 1{,}5~млн событий/с & Обработка трафика 10~Гбит/с при среднем размере пакета 1~КБ \
% Задержка диспетчеризации (p99) & ≤
% ≤ 100~мкс & Требования SLA для систем мониторинга реального времени \
% Потребление памяти на событие & ≤
% ≤ 256~Б & Эффективная работа с пулами при высокой нагрузке \
% Время горячей перезагрузки модуля & ≤
% ≤ 200~мс & Минимизация окон уязвимости при обновлении правил \
% Масштабируемость & Линейная до 32 потоков & Поддержка многоядерных серверных архитектур \
% \bottomrule
% \end{tabular}
% \end{table}

В качестве основного языка разработки выбран C++ (стандарт C++17/20), обеспечивающий статическую типобезопасность, zero-cost абстракции, детерминированное управление ресурсами и прямой доступ к аппаратным возможностям процессора.
Инструментарий включает систему сборки CMake для управления зависимостями и кроссплатформенной компиляции, Git для контроля версий, Clang/GCC в качестве компиляторов с поддержкой оптимизаций уровня -O3 и LTO, а также AddressSanitizer и ThreadSanitizer для выявления ошибок управления памятью и состояний гонки на ранних этапах разработки.
Выбор C++ обоснован его доминированием в высокопроизводительных сетевых системах (DPDK, Suricata, Zeek core) и наличием зрелой экосистемы шаблонного метапрограммирования, позволяющей реализовать событийную модель без runtime-накладок \cite{cpp2023perf}.

\subsection{Проектирование API системы}
В рамках проектирования API решаются три ключевые задачи: механизм создания новых типов событий, способ генерации (диспетчеризации) событий и метод регистрации обработчиков.
Остальные аспекты, включая управление потоками, синхронизацию и распределение памяти, инкапсулируются внутри библиотеки и прозрачны для пользователя.
Для следования лучшим практикам C++ и предотвращения конфликтов имён вся функциональность инкапсулирована в пространство имён \texttt{event}.

\subsubsection{Создание новых типов событий}
Для определения новых типов событий применяется механизм наследования на основе Curiously Recurring Template Pattern (CRTP).
Данный подход позволяет избежать виртуальных таблиц и гарантировать нулевые накладные расходы при приведении типов.
Пользователь определяет структуру события, наследуясь от шаблонного базового класса, который автоматически генерирует уникальный идентификатор типа и предоставляет методы доступа к полям.

\begin{lstlisting}[language=C++, caption={Определение пользовательского типа события}]
namespace event {
class IdGenerator {
 public:
  static size_t nextId() {
    static std::atomic<size_t> id_generator_{0};
    return id_generator_.fetch_add(1, std::memory_order::relaxed);
  }
};

template <typename Derived>
class EventBase : private IdGenerator {
 public:
  static size_t getTypeId() noexcept {
    static const size_t type_id = nextId();
    return type_id;
  }

 protected:
  ~EventBase() = default;
};
} // namespace event

struct PacketEvent : public event::BaseEvent<PacketEvent> {
    uint32_t protocol;
    uint16_t size;
    uint64_t timestamp_ns;
};
\end{lstlisting}

Использование CRTP в сочетании с \texttt{IdGenerator} гарантирует, что каждый тип события имеет уникальный линейный идентификатор.
Линейный рост идентификаторов позволяет эффективно адресоваться к массивам обработчиков за O(1), не требуя использования хеш-таблиц для маршрутизации.
Поля структуры располагаются в памяти непрерывно, что улучшает локальность кэша и позволяет применять векторизованные операции фильтрации при обработке групп событий.

% TODO: таким образом получаем гибкую систему настройки различных событий

\subsubsection{Генерация события заданного типа}
Для генерации событий реализуется шаблонная функция \texttt{dispatch}, принимающая произвольное число аргументов через \textit{perfect forwarding}.
Функция должна обеспечивать типобезопасность на этапе компиляции и эффективное создание объектов событий.

\begin{lstlisting}[language=C++, caption={Механизм диспетчеризации событий}]
template <typename EventType, typename... Args>
void dispatch(Args&&... args) {
    static_assert(std::is_base_of_v<event::BaseEvent<EventType>, EventType>,
                  "EventType must inherit from BaseEvent");

    // Event dispatch code
}
\end{lstlisting}

\subsubsection{Создание обработчиков событий фиксированного типа}
Регистрация обработчиков осуществляется через шаблонную функцию \texttt{subscribe}, принимающую любой вызываемый объект (лямбда, функция, функтор) с сигнатурой, совместимой с принимаемым типом события.

\begin{lstlisting}[language=C++, caption={Регистрация обработчика событий}]
template <typename EventType, typename Callable>
void subscribe(Callable&& callback) {
    static_assert(std::is_invocable_v<Callable, const EventType&>,
                  "Callback must accept const EventType&");

    // Handle subscribe code
}
\end{lstlisting}

\subsection{Внутренняя реализация и оптимизации}
\subsubsection{Type Erasure для обработчиков событий}
Для хранения разнородных обработчиков в едином контейнере применяется техника \textit{type erasure}, позволяющая привести произвольные вызываемые объекты к общему интерфейсу без потери типобезопасности на этапе вызова.
В отличие от стандартного \texttt{std::function}, в реализации используется реализации \texttt{fu2::unique\_function} из библиотеки \texttt{Naois/function2} \cite{function2_benchmark}.


\begin{lstlisting}[language=C++, caption={Handler с type erasure}]
namespace Core::event {

class Handler {
 public:
  template <typename EventType, typename Callable>
  Handler(std::type_identity<EventType>, Callable&& callback) {
    static_assert(
        std::is_invocable_v<Callable, const EventType&>, "Callback must accept const EventType&"
    );
    invoke_ = [cb = std::forward<Callable>(callback)](const void* event) {
      cb(*static_cast<const EventType*>(event));
    };
  }

  void handle(const void* event) {
    assert(invoke_);
    invoke_(event);
  }

 private:
  struct SBOCapacity {
    static constexpr std::size_t capacity  = 8 * 6;
    static constexpr std::size_t alignment = alignof(std::max_align_t);
  };

  template <typename... Signatures>
  using unique_function = fu2::function_base<true, false, SBOCapacity, true, false, Signatures...>;
  using InvokeFunction  = unique_function<void(const void*)>;

  InvokeFunction invoke_{};
};

}  // namespace Core::event
\end{lstlisting}

Выбор \texttt{fu2::function} вместо \texttt{std::function} обусловлен следующими факторами:
\begin{itemize}
    \item \textbf{Расширенный SBO}: стандартный \texttt{std::function} обычно имеет буфер 16--24 байта, что недостаточно для большинства лямбд с захватом нескольких переменных. Настройка буфера в 48 байт покрывает $\sim$95\% типичных обработчиков без динамической аллокации \cite{sbo_performance};
    \item \textbf{Больше возможностей}: поддержка только перемещаемых функций, поддержка \texttt{const}, \texttt{volatile}, \texttt{noexcept}.
\end{itemize}

\subsubsection{Реализация хранения и запуска обработчиков}
Мы будем хранить наши обработчики в такой структуре.
\begin{lstlisting}[language=C++, caption={}]
std::vector<std::vector<Handler>> handlers_storage;
\end{lstlisting}

В ячейке соответствующей идентификатору типа будет лежать массив с обработчиками событий этого типа.

Тогда функция \texttt{subscribe} будет просто сохранять обработчик в этот масив.

\begin{lstlisting}[language=C++, caption={Реализация подписки на событие}]
template <typename EventType, typename Callable>
void subscribe(Callable&& callback) {
  static_assert(std::is_invocable_v<Callable, const EventType&>,
                        "Callback must accept const EventType&");

  const std::size_t id = EventType::getTypeId();
  auto& handlers       = detail::handlerStorage().sync;

  if (id >= handlers.size()) {
    handlers.resize(id + 1);
  }

  handlers[id].push_back(
    Handler{std::type_identity<EventType>(), std::forward<Callable>(callback)}
  );
}
\end{lstlisting}

Получаем идентификатор типа события, на которое мы подписываемся, и хранилище обработчиков, расширяем хранилище если нужно и сохраняем обработчик.

\subsubsection{Оптимизация через синхронный запуск}

\newpage
% Архитектура библиотеки строится на принципах разделения ответственности: слой захвата (вне библиотеки), слой нормализации, ядро событийной шины, пулы обработки и слой доставки. Центральным компонентом выступает асинхронный маршрутизатор событий, реализованный на основе направленного ациклического графа (DAG) очередей. Каждый узел графа соответствует типу события или группе обработчиков, а рёбра определяют маршруты доставки с учётом приоритетов и фильтров.
%
% Для обеспечения высокой пропускной способности применяются следующие оптимизации:
% \begin{itemize}
%     \item Lock-free MPMC-очереди на основе атомарных индексов и cache-line выравнивания, исключающие false sharing \cite{wang2021wasm}.
%     \item Пулы памяти с pre-allocated блоками, специфичными для типов событий, что устраняет фрагментацию и ускоряет аллокацию/освобождение.
%     \item Thread-per-core модель с work-stealing для балансировки нагрузки при неравномерном распределении событий.
%     \item Векторизованная фильтрация событий на этапе маршрутизации с использованием SIMD-инструкций для параллельного сравнения полей.
%     \item Backpressure-механизм, адаптирующий скорость генерации к пропускной способности обработчиков через динамическое ограничение диспетчеризации и сброс низкоприоритетных событий при перегрузке \cite{dag2022design}.
% \end{itemize}
%
% Управление конфигурацией и схемами событий вынесено в декларативный реестр, поддерживающий версионирование и атомарные обновления. Это позволяет обновлять логику обработки без остановки конвейера, что критично для систем непрерывного мониторинга. Изоляция обработчиков достигается через разделение контекстов выполнения и ограничение доступа к глобальному состоянию, что предотвращает каскадные отказы при ошибках в пользовательской логике.
%
% \newpage
% TODO:
% Выделение памяти осуществляется из lock-free пула, специфичного для каждого типа события, что устраняет фрагментацию и contention на уровне менеджера памяти.
% Маршрутизация выполняется по типу через предвычисленные индексы, обеспечивая $O(1)$ сложность диспетчеризации.
%
% Для минимизации overhead применяется тип-стирание с опциональной оптимизацией: если обработчик не сохраняет состояние и совместим с указателем на функцию, используется прямой вызов, в противном случае применяется компактный \texttt{std::function}-аналог с small object optimization.
%
% Архитектура гарантирует потокобезопасность регистрации через lock-free атомарные операции и блокировки только на уровне метаданных реестра.
% Обработчики выполняются в контексте worker-потоков с привязкой к ядрам CPU (CPU affinity), что минимизирует переключения контекста и улучшает предсказуемость задержек. Система поддерживает множественных подписчиков на один тип события с настраиваемым порядком вызова (параллельным или последовательным), что позволяет реализовать конвейерную обработку без блокировок.
